\begin{abstract}
\noindent Quantum Relativity (QR) is developed as a theory of relational complexity governed by conservation of information. Registered relational structures undergo admissible transductions that preserve the distinctions needed to reconstruct the complete input. A proposed native functional $\mathcal C(R)$ measures relational complexity; $S(R)=G^{\mathcal C(R)}$, with fixed dimensionless $G>1$, is an optional scalar encoding. Phase and Mode describe organization and orientation, not interchangeable scalar amounts. Information conservation applies to the complete propagated-plus-retained output.

An execution is situated at a selected registration interface. Selecting a reference fixes presentation coordinates without making the reference physically privileged. Different situated accounts require orientation comparisons at their admissible incidences. This revision proves reference-change identities and a common-domain reconciliation criterion: for two sign-cocycle actions over a transitive incidence domain, an equivariant comparison exists exactly when their isotropy characters agree. It then has two globally sign-related representatives. Selecting a reference does not supply a missing common domain or repair incompatible transport.

The Thalean carrier $\mathrm{AT4val}[60,6]$ retains its recovery, receipt, history, measure and causal results, with $\cth=1$ under the declared locality and saturation hypotheses. The cover family, relational kernel, Mode action, history conference identity $K_{\rm hist}^2=5I$ and masses $w_\pm=(5\pm\sqrt5)/10$ remain unchanged. New finite replays reconstruct 120 native signed face incidences, 240 phase-decorated presentations, and transport of the order-ten realization kernels. They distinguish the split signed-axis group $A_5\times C_2$ from the analyzer-sector group $S_5$, and verify a literal registered-history phase action-groupoid automorphism. That phase theorem is not silently promoted to an intertwiner on the original two history labels.

The conservative-generator, gyroscopic and Maxwell calculations remain conventional correspondence models on specified domains. A native complexity functional, the common incidence and lifted actions coupling history to analyzers, the common-adjoint face and instrument realization, an endogenous nonperiodic history law, and controlled physical limits remain open. The contribution is a sharper architecture of situated registration and incidence-based comparison, not closure of those physical or native construction requirements.
\end{abstract}
\noindent\textbf{Keywords:} relational complexity; information-preserving transduction; selected reference; signed incidence; orientation reconciliation; history groupoid; conservative dynamics.
